\section{Background and motivation}
\label{sec:intro-background}

Autonomous robots that drive between crop rows underpin a range of
precision-agriculture tasks --- high-throughput phenotyping, targeted spraying,
mechanical weeding --- that larger over-canopy machinery cannot perform
\citep{sivakumar2021}. Each of them depends on the same primitive capability:
staying in the row, reliably, for the length of a season.

That capability is limited by localisation. Satellite positioning, the default
answer elsewhere in outdoor robotics, degrades under plant canopy;
\citet{higuti2019} report measuring over \SI{0.4}{\metre} of error beneath a
sorghum canopy with a high-end real-time kinematic receiver, and
\citet{navone2025} identify tall, dense vegetation as the primary cause of
satellite signal blockage in tree-row environments. Vision is the practical
alternative, and the question that follows is how the robot should represent
what it sees.

This dissertation sits at the intersection of computer vision, robotics and
control that defines the MSc in Robotics and Artificial Intelligence: a learned
perception module, a geometric stage that converts its output into a
navigable line, and a controller that converts that line into motion. The
contribution is not a new architecture for any one of those stages, but a
controlled measurement of how a design decision taken in the first propagates
into the last.

Three terms are used throughout in their established senses
\citep{kirillov2019}. \emph{Semantic} segmentation assigns a class label to
every pixel without separating individual objects. \emph{Instance}
segmentation detects and segments each object separately. \emph{Panoptic}
segmentation combines the two. The distinction matters here because the
formulations compared in this study do not all belong to the same paradigm.

\section{The problem}
\label{sec:intro-problem}

Vision-based row following has not converged on a single representation, but
one recurring choice is a binary crop mask: every pixel assigned either to
vegetation or to background, with all structure collapsed into a single
foreground class. \citet{desilva2024} state the design decision plainly,
representing a crop row ``as a single object rather than individual plants'',
and \citet{aghi2021} carry the same representation into vineyards.

The formulation is under challenge, though not from a single direction.
Recent work objects that it depends on the sky being visible above the row
\citep{navone2025}, that the free-space region it relies on degenerates under
sparse vegetation \citep{liu2023}, that pixel-wise segmentation is the wrong
intermediate representation altogether \citep{sivakumar2024}, and that learned
segmentation costs more annotation and computation than it returns
\citep{ahmadi2022}. What none of these objections addresses is the question of
how many classes such a mask should carry.

That question is worth asking in a vineyard because a vineyard row is not a
continuous band of foliage. It is defined by discrete rigid elements --- vine
trunks rooted in the ground and vertical trellis poles --- which remain
geometrically meaningful whether or not the sky is visible and whether or not
foliage has closed over the corridor. Class identity is already exploited in
vineyards on this basis, but for localisation rather than navigation, and the
two published class selections disagree with one another: \citet{papadimitriou2022}
retain trunks for their stability of appearance across seasons, while
\citet{desilva2025} retain trunks and discard poles because uniform shapes
frustrate feature matching. Both criteria belong to feature matching. Neither
asks what a line-fitting stage needs.

The question has therefore not been tested where it would bite: downstream of
perception, in a crop whose structure survives canopy closure, and under
conditions where the answer can be attributed to the formulation rather than
to the reference it is measured against.

\section{Research question and scope}
\label{sec:intro-question}

The study asks whether the choice of segmentation formulation --- a binary
crop mask, or a class-aware representation preserving trunk and pole identity
--- materially affects downstream centreline estimation and steering-command
generation for in-row vineyard navigation, and how large any such effect is.

The question is posed as a measurement rather than as a directional
hypothesis. A finding that formulation choice does not affect the navigation
output is as useful to a practitioner as a finding that it does, provided the
comparison is controlled tightly enough, and the resulting interval narrow
enough, that indistinguishability cannot be mistaken for insufficient
measurement.

Three perception arms share a single downstream pipeline.The first is a semantic model producing a binary crop mask, in which all structure is collapsed into a single foreground class. The second is an instance model producing the same binary
mask. The third is
an instance model preserving trunk and pole as distinct classes. Only the
second-to-third contrast is controlled: it holds architecture, training regime
and data fixed and varies the label scheme alone. The first-to-second contrast
varies many things at once and is treated throughout as a comparison of
complete systems rather than an attribution to architecture.

Three constraints bound the work. The evaluation is offline, against recorded
robot deployments, which supports per-frame analysis but not closed-loop
tracking. The reference trajectory in those recordings was produced under
autonomous topological navigation \citep{polvara2024} rather than by a human
driver, so the executed steering reflects a different sensing
modality from visual row following and cannot serve as a tuning target for the
visual controller; controller gains are instead derived from first principles.
And the annotation scheme folds foliage into background, so foliage cannot be
exploited as a navigation cue.

\section{Contributions}
\label{sec:intro-contributions}

This dissertation makes four contributions.

First, a controlled comparison of binary against class-aware segmentation for
downstream centreline extraction in a discrete-structure crop. No such
comparison appears in the vineyard navigation literature reviewed here, and
the controlled contrast is verified empirically rather than asserted.

Second, a null result established at two independent levels of measurement.
The three arms produce centreline estimates that are statistically
indistinguishable on the primary geometric metric across four recordings
spanning bare-vine and canopy conditions, and steering commands that are
likewise indistinguishable. The intervals are narrow enough to bound any
effect to a small fraction of the reference's own precision, so the result is
not simply an absence of evidence; and the two levels of measurement, which
fail in different ways, agree.

Third, an evaluation of availability alongside accuracy. The multiclass arm
yields a usable centreline less often than either binary arm, consistently in
direction across every dataset evaluated, including recordings from a later
season and from a second site with different camera hardware. Measuring both
outcomes gives a practitioner both halves of an answer that either alone would
misrepresent: the additional class detail buys no measurable accuracy and
costs availability.

Fourth, an account of what the comparison depends on. The error measured
against the recorded trajectory is separated into a component contributed by
the reference and a component contributed by the pipeline. The reference
component is shown to affect all three arms identically, and the arms are
therefore compared only through paired differences, in which it cancels. Six
independent groups report that no absolute reference is obtainable in this
setting; this study establishes rather than assumes that its findings do not
rest on one.